\documentclass[10pt,a4paper]{amsart}
\usepackage[margin=19mm]{geometry}
\usepackage{amsmath,amssymb,mathtools}
\usepackage[scr=rsfs,cal=euler]{mathalfa}
\usepackage{xfrac}
\usepackage[unicode,hidelinks,bookmarks=false]{hyperref}
\newcommand{\ie}{i.e.\ }
\newcommand{\cf}{cf.\ }
\newcommand{\resp}{respectively\ }
\newcommand{\Subsection}{\subsection}
\providecommand{\nobreakdash}{\nobreak\hbox{-}}
\setlength{\emergencystretch}{2em}
\multlinegap0pt
\usepackage{amssymb}
\usepackage{xcolor}
\usepackage{float}
\usepackage{mathtools}
\usepackage{stackengine}
\usepackage[shortlabels]{enumitem}
\newtheorem{prop}{Proposition}
\newtheorem{lem}{Lemma}
\title{On Hardy-Littlewood-Sobolev estimates for degenerate Laplacians}
\author{Pascal Auscher, Khalid Baadi}
\date{Source: arXiv:2506.20368v3 (CC BY 4.0)}
\begin{document}
\maketitle

\noindent\emph{Source and licence.} On Hardy-Littlewood-Sobolev estimates for degenerate Laplacians. Version arXiv:2506.20368v3 (CC BY 4.0), distributed by arXiv under the Creative Commons Attribution 4.0 International licence, http://creativecommons.org/licenses/by/4.0/, as recorded in the arXiv licence metadata for this exact version and shown on the abstract page https://arxiv.org/abs/2506.20368v3. Full licence notice: \texttt{licenses/arxiv-2506.20368-CC-BY-4.0.txt}.\par\smallskip

\noindent\textit{Editorial scope.} Counted unit: the article's lem lem:HLS (source0.tex lines 403--410). The article's complete original proof of it is delivered with it (source0.tex lines 411--456, one \texttt{proof} environment of 46 lines). No mathematics is written, repaired or re-derived: the statement and the proof are the article's own lines, in the article's own notation, and no retained line is substituted or re-flowed.  Retained uncounted as the source-local context the proof needs: lines 269--281; lines 309--311; lines 338--349.  Nothing was omitted from any retained span, so the package carries no omission record.  No retained line contains a citation, so the wrapper carries no bibliography and no bibliography entry.  No retained line contains a figure environment, an image-inclusion command or drawing code, so no image is delivered and the package declares no asset dependency.  Editorial formatting only, in addition: an AMS \texttt{amsart} preamble that restates the article's own declarations and macros used by the retained lines, namely the environments \texttt{prop}, \texttt{lem} on separate counters, together with the standard packages those lines need.  The retained environments print in this document as Proposition 1, Lemma 1 in order of appearance, whereas the article prints the same items with the numbering of its own full document.  The complete native source of the article as distributed in the arXiv e-print for this exact version is bundled unchanged as \texttt{sources/180491/source0.tex}.
\begin{prop}\label{prop:weights} We have the following properties.
    \begin{enumerate}
        \item $A_p(\mathbb{R}^n) \subset A_q(\mathbb{R}^n)$ for all $1 \leq p \leq q \leq \infty$.
        \item $RH_q(\mathbb{R}^n) \subset RH_p(\mathbb{R}^n)$ for all $1 < p \leq q \leq \infty$.
        \item If $\omega \in A_p(\mathbb{R}^n)$, $1<p<\infty$, then there exists $q \in (1,p)$ such that $\omega \in A_q(\mathbb{R}^n)$.
        \item If $\omega \in RH_q(\mathbb{R}^n)$, $1<q<\infty$, then there exists $p \in (q,\infty)$ such that $\omega \in RH_p(\mathbb{R}^n)$.
        \item $A_\infty(\mathbb{R}^n)= \bigcup_{1\le p<\infty} A_p(\mathbb{R}^n)= \bigcup_{1< q\le \infty} RH_q(\mathbb{R}^n) $.
        \item If $1\leq p \leq \infty$ and $1<q<\infty$, then $\omega \in A_p(\mathbb{R}^n) \cap RH_q(\mathbb{R}^n)$ if and only if $\omega^q \in A_{q(p-1)+1}(\mathbb{R}^n)$.
        \item If $1 \le p \le q < \infty$, then $\omega \in A_{p,q}(\mathbb{R}^n)$ if and only if $\omega^q \in A_{1+q/p'}(\mathbb{R}^n)$ if and only if $\omega \in A_{1+1/p'}(\mathbb{R}^n)\cap RH_q(\mathbb{R}^n)$.
        \item If $1 \le  p < q < \infty$, then $\omega^{1/p} \in A_{p,q}(\mathbb{R}^n)$ if and only if 
        $\omega \in A_p(\mathbb{R}^n) \cap RH_{q/p}(\mathbb{R}^n)$.
    \end{enumerate}
\end{prop}

\begin{equation}\label{DoublingMuck}
    \omega(2Q)\leq D \omega(Q),
\end{equation}

\begin{lem}\label{lem:GULB}
For all $t>0$ and $f \in L^2_\omega(\mathbb{R}^n)$, we have for almost every $x\in \mathbb{R}^n$,
    \begin{equation*}
        e^{t\Delta_\omega}f(x)= \int_{\mathbb{R}^n} K_t(x,y)f(y) \, \mathrm{d}\omega(y),
    \end{equation*}
    where $K_t(x,y)=\Gamma(t,x;0,y)$ and $\Gamma(t,x;s,y)$ is the generalized fundamental solution for the weighted heat operator $\partial_t-\Delta_\omega$. Moreover, there exist two constants $1<C=C([ \omega  ]_{A_2},n)<\infty$ and $1<c=c([ \omega  ]_{A_2},n)<\infty$ such that 
\begin{equation}\label{eq:GULB}
     \frac{C^{-1}}{\min(\omega_{t}(x),\omega_t(y))} e^{- \frac{c\left | x-y \right |^2}{t}} \leq 
     K_t(x,y) \leq \frac{C}{\max(\omega_{t}(x),\omega_t(y))} e^{- \frac{\left | x-y \right |^2}{ct}},
\end{equation}
for all $t>0$ and for (almost) all $(x,y) \in \mathbb{R}^{2n}$. 
\end{lem}

\begin{lem}\label{lem:HLS}
    Let $1 < p < q<\infty$. Assume that $\omega \in RH_{\frac{q}{p}}(\mathbb{R}^n)$. Then, there exist two constants $C = C([ \omega  ]_{A_2},[ \omega  ]_{RH_{\frac{q}{p}}},n,p,q)<\infty$ and $\varepsilon_{\omega, p,q}=\varepsilon(  [ \omega  ]_{RH_{\frac{q}{p}}},p,q,n)>0$ such that 
    \begin{equation*}
        \| e^{t\Delta_\omega} f \|_{L^q_{{\omega}^{q/p+\varepsilon}}(\mathbb{R}^n)} \leq C t^{-\frac{n}{2}(\frac{1}{p}-\frac{1}{q})} \|  f \|_{L^p_{\omega^{1+(p/q)\varepsilon}}(\mathbb{R}^n)},
        \end{equation*}
    for all $t>0$, $\varepsilon \in (-\varepsilon_{\omega, p,q},\varepsilon_{\omega, p,q})$ and all $f\in L^2_\omega(\mathbb{R}^n) \cap  L^p_{\omega^{1+(p/q)\varepsilon}}(\mathbb{R}^n)$. In particular, 
    $e^{t\Delta_\omega}$ extends to a bounded operator from $L^p_{\omega^{1+(p/q)\varepsilon}}(\mathbb{R}^n)$ to $L^q_{{\omega}^{q/p+\varepsilon}}(\mathbb{R}^n)$.
    \end{lem}

\begin{proof} The last point follows by density once the inequality is established and we prove it next. 

    Given that $\omega \in RH_{\frac{q}{p}}(\mathbb{R}^n)$, point (4) of Proposition \ref{prop:weights} and Hölder's inequality imply that there exist {$\varepsilon_{\omega, p,q} = \varepsilon([\omega]_{RH_{\frac{q}{p}}},p,q,n) \in \big(0, \min\big(\frac q p, \frac{q}{p'}\big)\big)$} and a constant $C = C([\omega]_{RH_{\frac{q}{p}}},p,q,n) <\infty$ such that, for all $\varepsilon \in (-\varepsilon_{\omega, p,q}, \varepsilon_{\omega, p,q})$ and all cubes $Q \subset \mathbb{R}^n$, the following inequalities hold
    \begin{equation}\label{eq:theta}
        \left ( \frac{1}{|Q|} \int_Q \omega^{\frac{q}{p}+\varepsilon}(x) \ \mathrm{d}x   \right )^{\frac{1}{\frac{q}{p}+\varepsilon}} \leq \frac{C}{|Q|} \int_Q \omega(x) \ \mathrm{d}x,
    \end{equation} 
    \begin{equation}\label{eq:thetatilde}
        \left ( \frac{1}{|Q|} \int_Q \omega^{1-\frac{p'}{q}\varepsilon}(x) \ \mathrm{d}x   \right )^{\frac{1}{1-\frac{p'}{q}\varepsilon} }\leq  \frac{C}{|Q|} \int_Q \omega(x) \ \mathrm{d}x.
    \end{equation}
    We fix $\varepsilon \in (-\varepsilon_{\omega, p,q},\varepsilon_{\omega, p,q})$ and set $\gamma= \frac{1}{p}+\frac{\varepsilon}{q} \in (0,1)$. We fix $f\in L^p_{\omega^{1+(p/q)\varepsilon}}(\mathbb{R}^n) \cap L^2_\omega(\mathbb{R}^n)$ . We fix $t>0$ and let $m \in \mathbb{Z}$ be the integer verifying $\frac{\sqrt{t}}{2}<\frac{1}{2^m}\le \sqrt{t}$. Let $(Q_{k})_{k\in \mathbb{Z}^n}$ the family of dyadic cubes with side lengths $\frac{1}{2^m}$. Using the upper bound for the heat kernel in Lemma \ref{lem:GULB}, we have  
    \begin{align*}
        \| e^{t\Delta_\omega} f &\|^q_{L^q_{{\omega}^{q/p+\varepsilon}}(\mathbb{R}^n)}= \int_{\mathbb{R}^n} |e^{t\Delta_\omega} f(x)|^q\  \omega^{q/p+\varepsilon}(x) \ \mathrm{d}x 
        \\&\leq C \int_{\mathbb{R}^n} \bigg( \int_{\mathbb{R}^n} \frac{e^{-\frac{|x-y|^2}{ct}}}{\omega_t(x)^{\gamma}\omega_t(y)^{1-\gamma}}|f(y)| \ \mathrm{d} \omega(y) \bigg)^q\  \omega^{q/p+\varepsilon}(x) \ \mathrm{d}x 
        \\& \leq C \sum_{k\in \mathbb{Z}^n} \int_{Q_{k}} \bigg( \sum_{\ell \in \mathbb{Z}^n} \int_{Q_{\ell}} \frac{e^{-\frac{|x-y|^2}{ct}}}{\omega_t(y)^{1-\gamma}}|f(y)| \frac{\omega^{\frac{\varepsilon}{q}}(y)}{\omega^{\frac{\varepsilon}{q}}(y)} \ \mathrm{d} \omega(y) \bigg)^q\ \frac{\omega^{q/p+\varepsilon}(x)}{\omega_t(x)^{q\gamma}} \ \mathrm{d}x
        \\& \leq C \sum_{k\in \mathbb{Z}^n} \int_{Q_{k}} \bigg( \sum_{\ell\in \mathbb{Z}^n} \bigg( \int_{Q_{\ell}} \frac{e^{-p'\frac{|x-y|^2}{ct}}}{\omega_t(y)^{p'(1-\gamma)}} \omega(y)^{-\frac{p'}{q}\varepsilon} \ \mathrm{d} \omega(y) \bigg)^{\frac{1}{p'}} \|\omega^{\frac{\varepsilon}{q}}  \, f \|_{L^p_\omega(Q_{\ell})} \bigg)^q \ \frac{\omega^{q/p+\varepsilon}(x)}{\omega_t(x)^{q \gamma}} \ \mathrm{d}x
        \\& \leq C \sum_{k\in \mathbb{Z}^n} \int_{Q_{k}} \bigg( \sum_{\ell \in \mathbb{Z}^n} e^{-\frac{d(Q_{k},Q_{\ell})^2}{ct}} \bigg( \int_{Q_{\ell}} \frac{\omega(y)^{1-\frac{p'}{q}\varepsilon}}{\omega_t(y)^{p'(1-\gamma)}} \ \mathrm{d} y \bigg )^{\frac{1}{p'}} \| f \|_{L^p_{\omega^{1+(p/q)\varepsilon}}(Q_{\ell})} \bigg)^q \ \frac{\omega^{q/p+\varepsilon}(x)}{\omega_t(x)^{q \gamma}} \ \mathrm{d}x 
        \\& = C \sum_{k\in \mathbb{Z}^n}  \bigg( \sum_{\ell \in \mathbb{Z}^n} e^{-\frac{d(Q_{k},Q_{\ell})^2}{ct}} \bigg( \int_{Q_{\ell}} \frac{\omega(y)^{1-\frac{p'}{q}\varepsilon}}{\omega_t(y)^{p'(1-\gamma)}} \ \mathrm{d} y \bigg)^{\frac{1}{p'}} \| f \|_{L^p_{\omega^{1+(p/q)\varepsilon}}(Q_{\ell})} \bigg)^q \int_{Q_{k}} \frac{\omega^{q/p+\varepsilon}(x)}{\omega_t(x)^{q \gamma}} \ \mathrm{d}x.
    \end{align*}
    Using the doubling property \eqref{DoublingMuck} with the equivalence of norms on $\mathbb{R}^n$, we see easily that there exists a constant $C=C(n,D)\ge 1 $ such that 
    $$ \frac{1}{C} \omega(Q_{k}) \leq \omega_t(x) \leq C \omega(Q_{k}), \quad \text{for all} \ k\in \mathbb{Z}^n \ \text{and all} \ x \in Q_{k}.$$
    Furthermore, for all $k,\ell \in \mathbb{Z}^n$ with $k \neq \ell$, we have
    $$d(Q_{k},Q_{\ell}) \ge C(n) \bigg(\frac{|k-\ell|}{2^m}-\frac{1}{2^m}\bigg) \ge C(n)(|k-\ell|-1)\frac{\sqrt{t}}{2}.$$
    Thus, we obtain 
    \begin{equation}\label{eq:smgroupe}
        \| e^{t\Delta_\omega} f \|^q_{L^q_{{\omega}^{q/p+\varepsilon}}(\mathbb{R}^n)} \leq C \sum_{k\in \mathbb{Z}^n}  \bigg( \sum_{\ell\in \mathbb{Z}^n} e^{-\frac{|k-\ell|^2}{c}} \frac{\big( \int_{Q_{\ell}} \omega(y)^{1-\frac{p'}{q}\varepsilon}\ \mathrm{d} y \big)^{\frac{1}{p'}}}{\omega(Q_\ell)^{1-\gamma}}  \| f \|_{L^p_{\omega^{1+(p/q)\varepsilon}}(Q_{\ell})} \bigg)^q \frac{\int_{Q_{k}} \omega^{\frac{q}{p}+\varepsilon}(x) \ \mathrm{d}x}{\omega(Q_k)^{q\gamma}}.
    \end{equation}
    In addition, applying inequality \eqref{eq:thetatilde} yields: 
    \begin{equation}\label{eq:1}
        \left ( \int_{Q_{\ell}} \omega(y)^{1-\frac{p'}{q}\varepsilon}\ \mathrm{d} y \right )^{\frac{1}{p'}} \leq  |Q_\ell|^{\frac{1}{p'}} \left( \frac{\omega(Q_\ell)}{|Q_\ell|} \right)^{\frac{1-\frac{p'}{q}\varepsilon}{p'}}=  |Q_\ell|^{\frac{\varepsilon}{q}} \omega(Q_\ell)^{\frac{1}{p'}-\frac{\varepsilon}{q}}= |Q_\ell|^{\frac{\varepsilon}{q}} \omega(Q_\ell)^{1-\gamma}.
    \end{equation}
    Likewise, since $\frac{q}{p} + \varepsilon = q \gamma$, it follows from \eqref{eq:theta} that:
    \begin{equation}\label{eq:2}
        \frac{\int_{Q_{k}} \omega^{q/p+\varepsilon}(x) \ \mathrm{d}x}{\omega(Q_k)^{q\gamma}} \leq C |Q_k|^{1-(\frac{q}{p}+\varepsilon)}.
    \end{equation}
    Since $|Q_k| = |Q_\ell| \sim t^{n/2}$, combining \eqref{eq:1} with \eqref{eq:2}, inequality \eqref{eq:smgroupe} takes the form:
    \begin{align*}
        \| e^{t\Delta_\omega} f \|^q_{L^q_{{\omega}^{q/p+\varepsilon}}(\mathbb{R}^n)} \leq &C t^{\frac{n}{2} \varepsilon} \times t^{\frac{n}{2}(1-(\frac{q}{p}+\varepsilon))} \sum_{k\in \mathbb{Z}^n}  \bigg ( \sum_{\ell\in \mathbb{Z}^n} e^{-\frac{|k-\ell|^2}{c}}   \| f \|_{L^p_{\omega^{1+(p/q)\varepsilon}}(Q_{\ell})} \bigg )^q 
        \\& = C \times t^{\frac{n}{2}(1-\frac{q}{p})} \sum_{k\in \mathbb{Z}^n}  \bigg ( \sum_{\ell\in \mathbb{Z}^n} e^{-\frac{|k-\ell|^2}{c}}   \| f \|_{L^p_{\omega^{1+(p/q)\varepsilon}}(Q_{\ell})} \bigg )^q. 
    \end{align*}
    Using Young's discrete convolution inequality with $s>1$ such that $\frac{1}{p}+\frac{1}{s}=1+\frac{1}{q}$, we deduce that 
    \begin{align*}
        \| e^{t\Delta_\omega} f \|^q_{L^q_{{\omega}^{q/p+\varepsilon}}(\mathbb{R}^n)} &\leq C t^{\frac{n}{2}(1-\frac{q}{p})} \bigg ( \sum_{k\in \mathbb{Z}^n} e^{-s\frac{|k|^2}{c}} \bigg )^{q/s} \bigg ( \sum_{k\in \mathbb{Z}^n}  \|f \|^p_{L^p_{\omega^{1+(p/q)\varepsilon}}(Q_{k})} \bigg )^{q/p} \\&= C t^{\frac{n}{2}(1-\frac{q}{p})} \bigg( \sum_{k\in \mathbb{Z}^n} e^{-s\frac{|k|^2}{c}} \bigg )^{q/s} \|f \|^q_{L^p_{\omega^{1+(p/q)\varepsilon}}(\mathbb{R}^n)}.
    \end{align*}
     Thus, we obtain
    $$\| e^{t\Delta_\omega} f \|_{L^q_{{\omega}^{q/p+\varepsilon}}(\mathbb{R}^n)} \leq C t^{-\frac{n}{2}(\frac{1}{p}-\frac{1}{q})} \|  f \|_{L^p_{\omega^{1+(p/q)\varepsilon}}(\mathbb{R}^n)}.$$
\end{proof}

\end{document}
